Pretrained 3-D CT foundation models are usually evaluated with a frozen encoder and a linear or attentive head. However, a group adapting such a model to new labels has to decide what to train, from a frozen probe through low-rank adapters to full fine-tuning, and whether to train one model per label or one for all labels of a cohort, and the evidence for these choices comes mostly from 2-D backbones and single tasks. In this work we compare \numClsMethodsN{} adaptation methods for multi-label classification and \numSegMethodsN{} for segmentation on four public 3-D CT encoders of three architectural families (TAP-CT, SPECTRE, Pillar-0 and VoCo) under one protocol: \numLabelsN{} binary labels from five public cohorts and \numStructuresN{} structures from five public datasets, each trained per task and jointly with one default setting and budget (\numRunsMainN{} runs, plus seed repeats and controls), and compared at label level with intervals over cohorts, labels and held-out cases. For classification, adapting the attention weights of the encoder improves on both the linear and the attentive frozen probe on the three main backbones, and one joint model per cohort is as accurate as one model per label; on the hierarchical VoCo encoder, adapting the encoder was level with the matched attentive probe, and both were behind its linear probe. For segmentation, a multi-depth decoder on the frozen encoder gives most of the gain over a linear probe, adapting the encoder on top of it changes Dice little on average, and one model per structure is more accurate than one joint model at our budgets. Full fine-tuning with its default setting failed to train on one backbone and recovered with layer-wise learning-rate decay, and the LoRA variants did not improve on plain LoRA. We condense the results into a practical guide that separates consistent findings from conditional and uncertain ones, and release code, splits and per-case predictions.
